\subsubsection{文脈の中のアーキテクチャ}
アーキテクチャは真空中で作られるものではない。どの開発ライフサイクルを使うか、製品系列の中にあるか、企業全体のアーキテクチャに従うか、システム・オブ・システムズの一部かによって、作り方が変わる。

伝統的なライフサイクルでは、要求に基づいてアーキテクチャ設計段階を置くことが多い。このとき、すべての要求が同じ重要度を持つわけではない。アーキテクチャに強く影響する要求、すなわち ASR を見極める必要がある。

製品系列や製品ファミリでは、共通のアーキテクチャが先に作られ、それを土台として個別製品が作られる。個別製品は、共通部分と可変部分を持つ。共通アーキテクチャは再利用性を高めるが、過度に固定すると個別製品の自由度を下げる。

アジャイル開発では、明示的なアーキテクチャ設計段階が小さいことがあり、コードからアーキテクチャが「創発する」と説明されることもある。この考え方は、利用者中心の情報システムではうまく働く場合がある。しかし、組込みシステム、サイバーフィジカルシステム、安全性が重要なシステムでは、利用者物語に現れにくい重要品質があるため、意識的なアーキテクチャ検討が必要である。

企業アーキテクチャやシステム・オブ・システムズの文脈では、上位のアーキテクチャが個別ソフトウェアの要求や制約になる。仕様、API、準拠試験、標準インタフェースなどを通じて制約が課される。
